% Source PDF: source/普通高考/2013/2013四川理.pdf, page 2 (question 18).
% Node -> directed-edge list: start -> input x -> x even?;
% no -> y=1 -> merge; yes -> x divisible by 3?;
% second yes -> y=3 -> merge; second no -> y=2 -> merge;
% merge -> output y -> finish. The three incoming merge arrows and the
% independent downward arrow into output are all preserved from the source.
\begin{tikzpicture}[
  >=Stealth,
  line width=0.45pt,
  line cap=round,
  line join=round,
  font=\normalsize,
  flow/.style={draw,anchor=center,align=center,inner sep=0pt,
    execute at begin node={\setlength{\parindent}{0pt}}},
  every node/.style={anchor=center,align=center,
    execute at begin node={\setlength{\parindent}{0pt}}},
  terminal/.style={flow,rounded corners=0.18cm,minimum width=1.45cm,minimum height=0.70cm},
  io/.style={flow,shape=trapezium,trapezium left angle=72,trapezium right angle=108,
    minimum width=2.20cm,minimum height=0.76cm},
  process/.style={flow,minimum height=0.72cm,minimum width=2.05cm},
  decision/.style={flow,diamond,aspect=2.6,inner sep=1pt,minimum width=4.25cm,minimum height=1.00cm},
  branchlabel/.style={inner sep=0pt,anchor=center,align=center,
    execute at begin node={\setlength{\parindent}{0pt}}},
]
  \node[terminal] (start) at (0,11.15) {开始};
  \node[io] (input) at (0,9.75) {输入 $x$};
  \node[decision,minimum width=4.55cm] (parity) at (0,8.15) {$x\text{ 为偶数?}$};
  \node[process,minimum width=1.85cm] (yone) at (4.05,8.15) {$y=1$};
  \node[decision,minimum width=5.05cm] (div3) at (0,6.45) {$x\text{ 能被 3 整除?}$};
  \node[process,minimum width=1.85cm] (ytwo) at (-4.05,6.45) {$y=2$};
  \node[process,minimum width=1.85cm] (ythree) at (0,5.00) {$y=3$};
  \coordinate (merge) at (0,3.95);
  \node[io,minimum width=2.30cm] (output) at (0,3.05) {输出 $y$};
  \node[terminal] (finish) at (0,1.55) {结束};

  \draw[->] (start.south) -- (input.north);
  \draw[->] (input.south) -- (parity.north);
  \draw[->] (parity.east) -- (4.05,8.15) -- (yone.west);
  \draw[->] (parity.south) -- (div3.north);
  \draw[->] (div3.west) -- (-4.05,6.45) -- (ytwo.east);
  \draw[->] (div3.south) -- (ythree.north);

  % Three branch outputs converge at the source's explicit merge point.
  \draw (yone.south) -- (4.05,3.95) -- (0.35,3.95) -- (merge);
  \draw (ytwo.south) -- (-4.05,3.95) -- (-0.35,3.95) -- (merge);
  \draw (ythree.south) -- (0,4.30) -- (merge);
  \draw[->] (merge) -- (output.north);
  \draw[->] (output.south) -- (finish.north);

  \node[branchlabel,anchor=south] at (3.35,8.78) {否};
  \node[branchlabel,anchor=west] at (0.20,7.46) {是};
  \node[branchlabel,anchor=east] at (-3.35,6.95) {否};
  \node[branchlabel,anchor=west] at (0.20,5.75) {是};
\end{tikzpicture}
